%% shortname: Existence with Production, Price Equilibrium with Transfers and the First Welfare Theorem
\documentclass[11pt]{article}
\usepackage{microstyle}

\begin{document}

\lessonheader{24}{Chapter 17 (end): existence of equilibrium in a production economy.
Chapter 16 (16.A--16.C): price equilibrium with transfers and the first welfare theorem.
An example of non-existence}{}


\section{Existence of equilibrium with production}

Finishing Chapter 17 today: the production economy. We need to modify our aggregate excess
demand function to account for production.

\topic{Assumptions on the production sets}
To have a more well-behaved problem, assume that the production sets are \emph{closed} and
\emph{bounded above} $\Rightarrow$ there is a solution to each firm's profit-maximization
problem. Assume also that each production set is \emph{strictly convex}
$\Rightarrow$ the solution is unique, so we have a supply \emph{function} rather than a
correspondence.

\topic{Profit and supply functions}
Let $\pi_j(p)$ denote firm $j$'s profit function and $y_j(p)$ its supply function:
\[ \pi_j(p)=p\cdot y_j(p), \qquad j=1,\dots,J. \]

\topic{Production-inclusive excess demand}
Define the production-inclusive aggregate excess demand function
\[ \tilde z(p)\;=\;\sum_{i=1}^{I}x_i\Bigl(p,\;p\cdot\omega_i+\sum_{j=1}^{J}\theta_{ij}\,\pi_j(p)\Bigr)
   \;-\;\sum_{i=1}^{I}\omega_i\;-\;\sum_{j=1}^{J}y_j(p). \]
Consumer $i$'s wealth is the value of the endowment plus the share $\theta_{ij}$ of the profit
of every firm $j$.

The function $\tilde z(p)$ satisfies the same five properties as $z(p)$ does (continuity,
homogeneity of degree zero, Walras' law $p\cdot\tilde z(p)=0$, bounded below, unbounded
above near the boundary), which we used to prove existence in a pure exchange economy
(Lesson 23). Hence
\begin{center}
\fbox{\parbox{0.86\linewidth}{\textbf{Result.} $p\gg 0$ is an equilibrium price vector if and
only if $\tilde z(p)=0$, and such a $p$ exists (same fixed-point argument as for $z(p)$).}}
\end{center}

\runin{Trick:} use Walras' law to cut down the degrees of freedom by 1 $\Rightarrow$ only
$L-1$ degrees of freedom, so we can normalize the price of one of the goods to be $1$.

\runin{Important:} prices are endogenous, and so is income, because income relies on selling
endowments and earning dividends.

\topic{Remarks}
\begin{flatlist}
  \item A solution exists under these assumptions.
  \item Convexity of the production sets is important. Even if it does not hold for
        individual firms, it may still hold in an approximate sense in the aggregate.
\end{flatlist}


\section{General equilibrium and welfare (MWG Chapter 16)}

\begin{flatlist}
  \item 16.A--16.C: the first welfare theorem.
  \item 16.D--16.F: the second welfare theorem, plus discussion.
\end{flatlist}

We wish to study the properties of Walrasian equilibrium, assuming it exists (sufficient
conditions for existence: last lecture).

\topic{Price equilibrium with transfers}
In fact, we extend the equilibrium concept beyond Walrasian equilibrium: a \emph{price
equilibrium with transfers}, i.e.\ a Walrasian equilibrium with wealth transfers.

\begin{flatlist}
  \item Given $p$, total wealth is
        \[ p\cdot\bar\omega+\sum_{j=1}^{J}p\cdot y_j^*, \qquad
           \bar\omega=\sum_{i=1}^{I}\omega_i . \]
  \item This total wealth can be redistributed. Let $w_i$ denote consumer $i$'s allocated
        wealth, with $\sum_{i=1}^{I}w_i=p\cdot\bar\omega+\sum_{j=1}^{J}p\cdot y_j^*$.
        $\Rightarrow$ consumer $i$'s budget set is
        \[ \{x_i\in X_i:\; p\cdot x_i\le w_i\}. \]
\end{flatlist}

\topic{Definition (price equilibrium with transfers)}
An allocation $(x^*,y^*)$ and a price vector $p$ constitute a price equilibrium with
transfers if there is some wealth distribution $(w_1,\dots,w_I)$ as above such that agents
are maximizing and demand $=$ supply:
\begin{flatnum}
  \item \textbf{Firms:} $p\cdot y_j\le p\cdot y_j^*$ for all $y_j\in Y_j$, $j=1,\dots,J$.
  \item \textbf{Consumers:} $x_i^*$ is optimal for $\succsim_i$ in
        $\{x_i\in X_i:\,p\cdot x_i\le w_i\}$, $i=1,\dots,I$.
  \item \textbf{Market clearing:} $\displaystyle\sum_{i=1}^{I}x_i^*=\bar\omega+\sum_{j=1}^{J}y_j^*$.
\end{flatnum}

\begin{flatlist}
  \item Walrasian equilibrium was the special case where there is no redistribution, but each
        consumer's wealth is the value of his endowment plus dividends:
        $w_i=p\cdot\omega_i+\sum_{j}\theta_{ij}\,p\cdot y_j^*$.
  \item $\Rightarrow$ The set of price equilibria with transfers is potentially a large set.
\end{flatlist}

Assume preferences are locally non-satiated. This does not guarantee existence, but if an
equilibrium exists, then it is ``nice'':

\topic{First fundamental welfare theorem}
\begin{center}
\fbox{\parbox{0.86\linewidth}{\textbf{Theorem.} If preferences are locally non-satiated and
$(x^*,y^*,p)$ is a price equilibrium with transfers, then the allocation $(x^*,y^*)$ is Pareto
efficient.

\textbf{Corollary.} Walrasian equilibrium allocations are Pareto efficient.}}
\end{center}

\runin{Proof} (by contradiction). Assume $(x^*,y^*)$ is not Pareto efficient
$\Rightarrow$ there is some feasible allocation $(x,y)$ that is a Pareto improvement over
$(x^*,y^*)$: $x_i\succsim_i x_i^*$ for all $i$ and $x_i\succ_i x_i^*$ for at least one $i$.

\runin{Step 0.} Local non-satiation implies $p\cdot x_i^*=w_i$ for all $i$ (otherwise there
would be a strictly better bundle close to $x_i^*$ that is still affordable).

\runin{Step 1: consumers who are made strictly better off.}
\[ x_i\succ_i x_i^* \;\Longrightarrow\; p\cdot x_i>w_i=p\cdot x_i^* . \]
Otherwise $x_i$ is affordable and the consumer should buy it instead of $x_i^*$
(contradicting that $x_i^*$ is optimal).

\runin{Step 2: no consumer is worse off.}
\[ x_i\succsim_i x_i^* \;\Longrightarrow\; p\cdot x_i\ge w_i=p\cdot x_i^* . \]
Otherwise $x_i$ is affordable with money to spare, so (by local non-satiation) an even better
bundle is affordable, contradicting that $x_i^*$ is optimal.

\runin{Step 3: sum over all agents,} keeping in mind that the inequality is strict for at
least one agent:
\[ \sum_{i=1}^{I}p\cdot x_i\;>\;\sum_{i=1}^{I}w_i
   \;=\;p\cdot\bar\omega+\sum_{j=1}^{J}p\cdot y_j^*
   \;\ge\;p\cdot\bar\omega+\sum_{j=1}^{J}p\cdot y_j , \]
where the last inequality is profit maximization ($p\cdot y_j^*\ge p\cdot y_j$ for every
$y_j\in Y_j$).

\runin{Step 4: quantities.} The above was expressed in money. Now consider quantities:
$(x,y)$ must be feasible,
\[ \sum_{i=1}^{I}x_i=\bar\omega+\sum_{j=1}^{J}y_j , \]
so their values at prices $p$ are the same:
\[ p\cdot\sum_{i}x_i=p\cdot\bar\omega+p\cdot\sum_{j}y_j
   \quad\Longleftrightarrow\quad
   \sum_{i}p\cdot x_i=p\cdot\bar\omega+\sum_{j}p\cdot y_j . \]
This contradicts the earlier conclusion that
$\sum_i p\cdot x_i>p\cdot\bar\omega+\sum_j p\cdot y_j$. $\Rightarrow$ $(x^*,y^*)$ is Pareto
efficient. \hfill$\blacksquare$


\section{Example: why strict convexity and strong monotonicity matter for existence}

A pure exchange economy with two agents and two goods ($x_{ki}$: consumption of good $k$ by
agent $i$).

\begin{flatlist}
  \item \textbf{Agent 1:} $u_1(x_{11},x_{21})=\sqrt{x_{11}}+x_{21}$ $\Rightarrow$
        indifference curves $x_{21}=c-\sqrt{x_{11}}$, with unbounded slope as
        $x_{11}\to 0$.
  \item \textbf{Agent 2:} $u_2(x_{12},x_{22})=x_{12}$. Locally non-satiated, but neither
        strongly monotone nor strictly convex (agent 2 does not care about good 2).
  \item \textbf{Endowments:} agent 1 has all of good 2, agent 2 has all of good 1:
        $\omega_1=(0,\bar\omega_2)$, $\omega_2=(\bar\omega_1,0)$.
\end{flatlist}

\begin{center}
\begin{tikzpicture}[x=1.15cm,y=1.15cm]
  % box: omega1bar = 5, omega2bar = 3; endowment at (0,3)
  % agent 1: u = sqrt(x) + y; prices p1/p2 = 1/2 -> demand x11 = 1, x21 = 2.5
  \draw[line width=0.6pt] (0,0) rectangle (5,3);
  \node[ecnode,below left] at (0,0) {$O_1$};
  \node[ecnode,above right] at (5,3) {$O_2$};
  % Pareto set: top edge
  \draw[adm] (0,3) -- (5,3);
  \node[ecnode,above] at (2.5,3.05) {Pareto efficient allocations};
  % agent 1 indifference curves: through omega (c=3) and tangent at x1(p) (c=3.5)
  \draw[budget,domain=0:5,samples=120] plot (\x,{3-sqrt(\x)});
  \draw[budget,domain=0.25:5,samples=120] plot (\x,{3.5-sqrt(\x)});
  \draw[-{Stealth}] (1.5,1.15) -- (1.9,1.55);
  \node[ecnode,anchor=west] at (0.9,0.4) {agent 1's indiff.\ curves};
  % agent 2 indifference curves (vertical)
  \draw[dashed] (4.65,0) -- (4.65,3);
  \draw[dashed] (4.3,0) -- (4.3,3);
  \draw[-{Stealth}] (4.95,2.3) -- (4.4,2.3);
  \node[ecnode,anchor=west] at (5.1,2.45) {agent 2's indiff.};
  \node[ecnode,anchor=west] at (5.1,2.1) {curves (vertical)};
  % budget line through omega, slope -p1/p2 = -1/2
  \draw[thin] (0,3) -- (5,0.5);
  \node[ecnode,anchor=west] at (5.1,0.6) {budget line, slope $-p_1/p_2$};
  % endowment and agent 1's demand
  \node[bundle] at (0,3) {};
  \node[ecnode,above left] at (0,3) {$\omega$};
  \node[bundle] at (1,2.5) {};
  \node[ecnode,above right] at (1,2.55) {$x_1(p)$};
  % labels
  \node[ecnode,below] at (2.5,-0.15) {good 1 (agent 1)};
  \node[ecnode,rotate=90,above] at (-0.15,1.5) {good 2 (agent 1)};
\end{tikzpicture}
\end{center}

The Pareto efficient allocations are the top edge of the box: agent 1 gets all of good 2 (which
agent 2 does not value), and good 1 is split in any way.

\topic{There is no equilibrium}
\begin{flatlist}
  \item \textbf{Look for equilibrium prices $p\gg 0$.} Agent 2 keeps her endowment of good 1:
        she only cares about good 1, and her wealth $p_1\bar\omega_1$ buys exactly
        $\bar\omega_1$ units of it. But agent 1 is also demanding good 1: his wealth
        $p_2\bar\omega_2$ is positive and his marginal utility of good 1 is unbounded at
        $x_{11}=0$, so $x_{11}(p)>0$. Hence total demand for good 1 exceeds the supply
        (endowment):
        \[ z_1(p)=x_{11}(p)+\bar\omega_1-\bar\omega_1=x_{11}(p)>0
           \;\Rightarrow\; \text{not an equilibrium.} \]
  \item \textbf{If $p_1=0$ or $p_2=0$}, there is a free good $\Rightarrow$ unbounded demand
        ($p_1=0$: agent 2 demands unboundedly much of good 1; $p_2=0$: agent 1 demands
        unboundedly much of good 2):
        \[ z_1(p)>0 \quad\text{or}\quad z_2(p)>0 . \]
\end{flatlist}
$\Rightarrow$ No equilibrium exists.


\begin{transcriptnotes}
  \item The export has 6 pages; page 6 is blank apart from three lines at the top. No date
        is written, so the header date is left empty; the subtitle and the section headings
        are \textbf{added}.
  \item Page 1: ``Finishing 17 today. Production Economy''. ``modify over aggregate excess
        demand'' $\to$ \emph{our}. ``production sets are chosen and burned above'' is
        read as \emph{closed and bounded above} (the standard assumptions guaranteeing a
        solution to profit maximization). ``Let $\pi_i(P)$ denote firm $j$'s profit function
        and $X_i(P)$ denote its supply function; $\pi_i(P)=Py_i(P)$'' is \textbf{corrected}
        to $\pi_j$, $y_j$ and $\pi_j(p)=p\cdot y_j(p)$.
  \item Page 1: ``Define production function inclusive excess demand function'' $\to$
        \emph{production-inclusive}. In $\tilde z(p)$ the upper limits of both $j$-sums
        are written $I$ (or unclear); \textbf{corrected} to $J$. The sentence on
        consumer $i$'s wealth is \textbf{added}.
  \item Page 1: ``same five properties as $z(p)$ does and which we used to prove existance
        in a pure exchange economy''; the list of the five properties in parentheses is
        \textbf{added} from the Lesson 23 slides. ``$p\gg0$ is an equilibrium price vector
        iff $\tilde z(p)=0$'' is kept and boxed; the clause ``and such a $p$ exists'' is
        \textbf{added} (it is what ``a solution exists'' on page 2 refers to). ``walra's law
        to cut down the degree of freedom by 1'' $\to$ \emph{Walras' law \dots\ degrees}.
  \item Page 2: ``normalize the price of one of the goods tobe 1''. ``Prices are endogenous
        so is income because it relies on selling endowments and earning divident'' $\to$
        \emph{dividends}. ``A solution exists under one assumptions'' $\to$ \emph{these
        assumptions}. ``Convexity \dots\ even it does not hold for individual firms it still
        hold in an approximate sense on the aggregate'' $\to$ \emph{it may still hold in an
        approximate sense in the aggregate}.
  \item Page 2: ``Chapter 16: Geneal equilibrium and welfane''; ``16 A--C first welfare
        therom''; ``16 D--F [section, crossed out] second \dots\ + discussion'' $\to$
        \emph{second welfare theorem, plus discussion}. ``study properties of walrasian
        equilibrium assuming it exists $\to$ sufficient conditions yestoday'' $\to$
        \emph{sufficient conditions for existence: last lecture}.
  \item Page 2: ``extend the EQ concept beyound walrasian equ'' $\to$ \emph{beyond Walrasian
        equilibrium}. Total wealth $p\cdot\bar\omega+\sum_j p\cdot y_j^*$ is kept; the
        definition $\bar\omega=\sum_i\omega_i$ and the condition
        $\sum_i w_i=p\cdot\bar\omega+\sum_j p\cdot y_j^*$ are \textbf{added}. Budget set
        ``$Px_i\le w_i$'' written as a set.
  \item Page 3: ``We now look for a price-equilibrium with transfers $(x^*,y^*,p)$ if there is
        some wealth distribution that agents are maximizing and demand $=$ supply'' is turned
        into the formal \textbf{Definition}; the three numbered conditions are
        \textbf{added} (MWG Definition 16.B.4), written as in Lesson 22.
  \item Page 3: ``Walrasian equilibrium was the special case welfare there is no
        redistribution but the consumer's wealth is the value of his endowment $+$
        Divident'' $\to$ \emph{where there is no redistribution \dots\ dividends}; the
        formula for $w_i$ is \textbf{added}. ``potentilly a large set''.
  \item Page 3: ``locally non-satiated. Does not generate existance'' $\to$ \emph{does not
        guarantee existence}. ``The first [crossed-out word] welfare theorem assume colar
        [crossed out] Non-satiation and that $(x^*,y^*,P)$ is a price equb.\ with transfers,
        then $(x^*,y^*)$ is P.E.'' $\to$ boxed theorem; P.E.\ $=$ Pareto efficient.
        ``Collary: Walrasian EQ are P.E.'' $\to$ \emph{Corollary}. The precise meaning of a
        Pareto improvement is \textbf{added} to the proof.
  \item Page 3: ``Before we get to implication, local non-satiation implies $P\cdot X_i^*=w_i$
        for $\forall i$'' $\to$ Step 0; the reason in parentheses is \textbf{added}.
  \item Page 4: Step 1 is written $x_i\succsim_i x_i^*\Rightarrow Px_i>w_i$;
        \textbf{corrected} to the strict preference $x_i\succ_i x_i^*$ (``consumers that are
        made strictly better off''; with weak preference the strict inequality does not
        follow). ``Contradicting $x_i^*$ is optimum''. Step 2: ``an even better [bundle,
        crossed out] bundle is affordable''; the appeal to local non-satiation is
        \textbf{added}.
  \item Page 4: ``Sum over all agent, keeping in mind that the inequality is strict for at
        least one agent''; in the chain, the annotation ``profit max'' under the last
        $\ge$ is kept as the explanation after the display. The second sum on the page is
        written $\sum_{i=1}^{I}w_i$ (kept); the summation limits on $y_j$ are kept as $J$.
        ``Let us now consider quemtities: $(x,y)$ must be feasible''.
  \item Page 5: ``so their value is the same'' and the two equalities are kept; ``contradicts
        the earlier conclusion that $\sum Px_i>P\bar\omega+\sum P\cdot y_j$
        $\Rightarrow(x^*,y^*)$ is P.E.''
  \item Page 5: ``An example of a pure exchange to show why strict convexity $+$ strong
        monotonicity is important for existence'' is used as the section title. Agent 1's
        utility is written $v(x_{11},x_{12})=\sqrt{x_{11}}+x_{21}$; \textbf{corrected} to
        $u_1(x_{11},x_{21})$ (good $k$, agent $i$, as in Lessons 19 and 21). The annotation
        under $\sqrt{x_{11}}$ reads ``good agent'' (uncertain; probably ``good 1, agent 1'')
        and is dropped. ``indiff curve $x_2=c-\sqrt{x_{11}}$ with unbounded score as
        $x_{11}\to0$'' $\to$ $x_{21}=c-\sqrt{x_{11}}$, \emph{unbounded slope}. Agent 2:
        $V(x_{12},x_{22})=x_{12}$, ``locally non-satiated''; the remark that agent 2 is
        neither strongly monotone nor strictly convex is \textbf{added}. The endowment
        vectors are \textbf{added}.
  \item Page 5, figure: Edgeworth box redrawn to scale ($\bar\omega_1=5$, $\bar\omega_2=3$,
        $p_1/p_2=1/2$, so agent 1 demands $x_1(p)=(1,2.5)$, where his indifference curve
        is tangent to the budget line). Kept from the page: $\omega$
        at the top-left corner, the brace ``P.E.\ allocation'' along the top edge (drawn as
        a thick line), agent 1's indifference curves with the arrow of increasing utility,
        ``agent 2's indiff curve'' (vertical) with an arrow pointing left, ``Agent 1'',
        ``Agent 2''. The budget line, the demand point $x_1(p)$, the origins $O_1$, $O_2$
        and the axis labels are \textbf{added}; the sentence explaining the Pareto set is
        \textbf{added}.
  \item Page 5: ``Look for EQ prices $p\gg0$ $\Rightarrow$ Agent 2 keeps her endowment of
        good 1 $+$ agent 1 is also demanding good 1. Total demand $>$ supply endow.\ of good
        1, $Z(P)>0$ $\Rightarrow$ Not EQ''; $Z(P)$ \textbf{corrected} to $z_1(p)$; the
        reasons (agent 2's wealth, agent 1's unbounded marginal utility) are \textbf{added}.
  \item Page 6: ``If $P_1$ and $P_2$ is $=0$'' $\to$ \emph{if $p_1=0$ or $p_2=0$}.
        ``$z(p)>0$ or $z_2(p)>0$'' $\to$ $z_1(p)>0$ or $z_2(p)>0$; which agent's demand
        explodes in each case is \textbf{added}.
\end{transcriptnotes}

\end{document}
